\documentclass[11pt,a4paper]{article}

%% PACHETE MATEMATICE
\usepackage{amsmath,amssymb,amsfonts}
\usepackage{amsthm}
\usepackage{mathpazo}
\usepackage{eulervm}
\usepackage{enumerate}
\usepackage{color}
\usepackage{graphicx}
\usepackage[all]{xy}
\usepackage{xcolor}
\usepackage{framed}
\usepackage[unicode]{hyperref}
\setcounter{MaxMatrixCols}{20}
\usepackage{multicol}
%\usepackage[romanian]{babel}


%% ASPECT
\linespread{1.05}
\usepackage[T1]{fontenc}
\usepackage[utf8x]{inputenc}
\usepackage[margin=2cm]{geometry}
% \usepackage[color]{showkeys}
\usepackage{anyfontsize}
\usepackage{titlesec}
\usepackage{tocloft}
\renewcommand{\cfttoctitlefont}{\hfill\large\textsc}
\renewcommand{\cftaftertoctitle}{\hfill}
\titleformat*{\section}{\large\bfseries}

\usepackage{fancyhdr}
\pagestyle{fancy}
\rhead{\small \color{gray} Notițe de Adrian Manea}
\lhead{\small \color{gray} ETTI-AM1, 2017---2018, Consultație}


\usepackage[autostyle = false, style = german]{csquotes}
\usepackage[romanian]{babel}
\newcommand\qq{\enquote}

%% COMENZILE MELE
\renewcommand*{\proofname}{Dem.:}
\newcommand\bb{\mathbb}
\newcommand\dr{\mathrm}
\newcommand\kal{\mathcal}
\newcommand\fk{\mathfrak}
\newcommand\op{\oplus}
\newcommand\ot{\otimes}
\newcommand\ds{\displaystyle}
\newcommand\id{\indent}
\newcommand\nid{\noindent}
\newcommand\seq{\subseteq}
\newcommand\sneq{\subsetneq}
\newcommand\speq{\supseteq}
\newcommand\spneq{\supsetneq}
\newcommand\rar{\rightarrow}
\newcommand\Rar{\Rightarrow}
\newcommand\xrar{\xrightarrow}
\newcommand\lar{\leftarrow}
\newcommand\Lar{\Leftarrow}
\newcommand\xlar{\xleftarrow}
\newcommand\lrar{\leftrightarrow}
\newcommand\Lrar{\Leftrightarrow}
\newcommand\ideal{\trianglelefteq}
\newcommand\ai{\text{ a.\^{i}. }}
\newcommand\sk{(\textit{Schi\c{t}\u{a}}) }
\newcommand\rhup{\rightharpoonup}
\newcommand\rhdn{\rightharpoondown}
\newcommand\lhup{\leftharpoonup}
\newcommand\lhdn{\leftharpoondown}
\newcommand\vid{\emptyset}
\newcommand\wh{\widehat}
\newcommand\ol{\overline}
\newcommand\NN{\mathbb{N}}
\newcommand\RR{\mathbb{R}}
\newcommand\ZZ{\mathbb{Z}}
\newcommand\QQ{\mathbb{Q}}
\newcommand\CC{\mathbb{C}}

%% STILURILE MELE
\newtheoremstyle{dotless}{}{}{\itshape}{}{\bfseries}{:}{ }{}

\theoremstyle{dotless}
\newtheorem{theorem}{Teorem\u{a}}[section]
\newtheorem{proposition}{Propozi\c{t}ie}[section]
\newtheorem{lemma}{Lem\u{a}}[section]
\newtheorem{corollary}{Corolar}[section]
\newtheorem{exercise}{Exerci\c{t}iu}

\newtheoremstyle{dotlessdr}{}{}{}{}{\bfseries}{:}{ }{}
\theoremstyle{dotlessdr}
\newtheorem{example}{Exemplu}[section]
\newtheorem{definition}{Defini\c{t}ie}[section]
\newtheorem{remark}{Observa\c{t}ie}[section]




\begin{document}

\begin{center}
  {\large\textbf{Consultație \\
      Extreme și integrale }}
\end{center}

\vspace{2cm}

1. Să se determine maximul și minimul funcției:
\[
  f(x, y) = x^2 + y^2 - 3x - 2y + 1
\]
pe mulțimea $K = \{ (x, y) \in \RR^2 \mid x^2 + y^2 \leq 1 \}$.

\emph{Indicații:}
\begin{itemize}
\item mulțimea $K$ este compactă (un disc), funcția este continuă, deci
  este mărginită și își atinge marginile;
\item studiind pe interiorul discului ($x^2 + y^2 < 1$), nu avem soluții;
\item pe frontieră, avem legătura $g(x, y) = x^2 + y^2 - 1 = 0$. Folosind metoda
  multiplicatorilor lui Lagrange, găsim punctele:
  \[
    (x, y) \in \Big\{ \Big( -\frac{3\sqrt{13}}{13}, -\frac{2\sqrt{13}}{13} \Big),
    \Big(\frac{3\sqrt{13}}{13}, \frac{2\sqrt{13}}{13}\Big) \Big\}
  \]
\item pentru a găsi punctele de extrem, putem calcula $f$ în fiecare din aceste valori.
\end{itemize}

\vspace{1cm}

2. Calculați:
\begin{enumerate}[(a)]
\item $\ds\int_0^\infty \frac{dx}{1 + x^3}$;
\item $\ds\int_0^1 \ln^p \Big( \frac{1}{x} \Big), p > -1$.
\end{enumerate}

\emph{Indicații:}
\begin{enumerate}[(a)]
\item Vom face schimbarea de variabilă $x^3 = y$ și folosim formula pentru funcția beta:
  \[
    B(p, q) = \int_0^\infty \frac{y^{p-1}}{(1 + y)^{p+q}} dy.
  \]
  Rezultatul este: $\frac{1}{3} B\big( \frac{1}{3}, \frac{2}{3} \big)$.
\item Vom face schimbare de variabilă $\ln x = y$ și folosim formula pentru funcția gamma:
  \[
    \Gamma(\alpha) = \int_0^\infty x^{\alpha - 1} e^{-x} dx, \alpha > 0.
  \]
  Rezultatul este $\Gamma(p + 1)$.
\end{enumerate}

\vspace{1cm}

3. Să se calculeze volumul mulțimii $\Omega$ mărginită de suprafețele de ecuații:
\[
  \begin{cases}
    2x^2 + y^2 + z^2 &= 1 \\
    2x^2 + y^2 - z^2 &= 0 \\
    z \geq 0
  \end{cases}
\]

\emph{Indicații:} Eliminînd $z$ din ecuații, găsim elipsa $4x^2 + 2y^2 = 1$, care
este curba de intersecție a suprafețelor. De asemenea, avem și $z = \frac{1}{\sqrt{2}}$.
Așadar, putem proiecta mulțimea pe planul $XOY$ și găsim:
\[
  D = \{ (x, y) \in \RR^2 \mid 4x^2 + 2y^2 \leq 1 \}.
\]
Cu definiția, rezultă:
\[
  V(\Omega) = \iiint_\Omega dx dy dz = \iint_D dxdy \int_{\sqrt{2x^2 + y^2}}^{\sqrt{1 - 2x^2 - y^2}} dz.
\]

Trecînd la coordonate polare, găsim rezultatul $\frac{\pi}{3} (\sqrt{2} - 1)$.

\vspace{1cm}

4. Să se calculeze circulația cîmpului de vectori $\vec{V}$ de-a lungul curbei $\Gamma$, unde:
\[
  \vec{V} = -(x^2 + y^2) \vec{i} - (x^2 - y^2)\vec{j},
\]
iar curba $\Gamma = \Gamma_1 \cup \Gamma_2$ este:
\begin{align*}
  \Gamma_1 &= \{(x, y) \in \RR^2 \mid x^2 + y^2, y < 0 \} \\
  \Gamma_2 &= \{(x, y) \in \RR^2 \mid x^2 + y^2 - 2x = 0, y \geq 0 \}.
\end{align*}

\emph{Indicații:} Oricărui cîmp vectorial $\vec{V} = P\vec{i} + Q\vec{j} + R\vec{k}$ i se asociază
o 1-formă diferențial $\alpha = Pdx + Qdy + Rdz$. Atunci circulația cîmpului $\vec{V}$ de-a lungul
curbei $\Gamma$ devine integrala curbilinie de speța a doua $\int_\Gamma \alpha = \int_\Gamma \vec{V} \cdot d\vec{r}$.

Avem nevoie să parametrizăm curbele. Vom folosi coordonate polare și găsim:
\begin{align*}
  \Gamma_1 &: \begin{cases}
    x(t) &= 2\cos t \\
    y(t) &= 2\sin t
  \end{cases}, t \in [\pi, 2\pi) \\
  \Gamma_2 &: \begin{cases}
    x(t) &= 1 + \cos t \\
    y(t) &= \sin t
  \end{cases}, t \in [0, \pi].
\end{align*}

Acum putem aplica formula de calcul pentru integrala curbilinie de speța a doua:
\[
  \int_\Gamma Pdx + Qdy + Rdz = \int_a^b (P \circ \Gamma)(t) \cdot x'(t) + (Q\circ \Gamma)(t) \cdot y'(t) +%
  (R \circ \Gamma)(t) z'(t) dt, \quad t \in [a, b].
\]

\vspace{1cm}

5. Să se calculeze aria paraboloidului $z = x^2 + y^2, z \in [0, h], h \in \RR$.

\emph{Indicații:} Deoarece paraboloidul este dat în forma unei \emph{parametrizări carteziene}, folosim
formula corespunzătoare de la integrale de suprafață de speța întîi. Așadar, avem $z = f(x,y) = x^2 + y^2$ și
aria se calculează folosind formula:
\[
  \int_\Sigma F(x, y, z) d\sigma = \int_D F(x, y, f(x, y)) \cdot \sqrt{1 + p^2 + q^2} dxdy,
\]
unde $p = \frac{\partial z}{\partial x}$, iar $q = \frac{\partial z}{\partial y}$, iar $F = 1$ pentru
calculul ariei.

\vspace{1cm}

6. Să se calculeze fluxul cîmpului de vectori $\vec{V}$ prin suprafața $\Sigma$ pentru:
\[
  \vec{V} = y\vec{i} - x\vec{j} + z^2 \vec{k}, \quad
  \Sigma: z^2 = x^2 + y^2, z \in [0,1].
\]

\emph{Indicații:} Folosind formula de definiție pentru flux, avem:
\[
  F_\Sigma(\vec{V}) = \int_\Sigma \vec{V} \cdot \vec{n} d\sigma = %
  \iint_D (\vec{V} \circ \Sigma) \cdot %
  \Big( \frac{\partial \Phi}{\partial u} \times \frac{\partial \Phi}{\partial v} \Big) dudv,
\]
unde $\Phi = \Phi(u, v) : D \to \RR^3$ este o parametrizare a suprafeței $\Sigma$.

Putem folosi parametrizarea carteziană $z = \sqrt{x^2 + y^2}$ sau, echivalent:
\[
  \begin{cases}
    x(u, v) &= u \\
    y(u, v) &= v \\
    z(u, v) &= \sqrt{u^2 + v^2}
  \end{cases}
\]

Putem calcula acum componentele \emph{versorului} normal, apoi produsul scalar și, în fine,
integrala.

\vspace{1cm}

7. Să se calculeze integrala curbilinie $\ds\int_\Gamma \alpha$, unde:
\[
  \alpha = ydx + x^2dy,
\]
iar $\Gamma$ este pătratul cu vîrfurile $A(0, 0), B(2, 0), C(2,2), D(0, 2)$.

\emph{Indicație:} Se folosește formula Green-Riemann, unde compactul $K$ este interiorul pătratului.
Integrala devine:
\[
  \iint_K (1-2y)dxdy = \int_0^2 \int_0^2 (1 - 2y) dy = -4.
\]

\vspace{1cm}

8. Fie forma diferențială $\alpha = -\frac{y}{x^2 + y^2} dx + \frac{x}{x^2 + y^2} dy$.

Să se calculeze $\ds\int_\Gamma \alpha$, unde $\Gamma$ este cercul de centru 0 și rază $R > 0$.

\emph{Indicație:} În acest caz \textbf{NU} se poate aplica formula Green-Riemann, deoarece
$\alpha$ nu este definită în $(0, 0)$, care se găsește în interiorul compactului delimitat de $\Gamma$
(mai general, conform teoriei, $\alpha$ trebuie să fie de clasă $\kal{C}^1$ în interior).

Așadar, folosim definiția și parametrizăm cercul, cu coordonate polare, apoi folosim
definiția integralei curbilinii de speța a doua.

\vspace{1cm}

9. Să se calculeze circulația cîmpului de vectori $\vec{V}$ pe curba $\Gamma$ pentru:
\[
  \vec{V} = y^2\vec{i} + xy\vec{j}, \quad \Gamma = \{(x, y) \in \RR^2 \mid x^2 + y^2 = 1, y > 0 \}.
\]

\emph{Indicații:} Pe lîngă metoda de la exercițiul 4, putem aplica și formula Green-Riemann.
Obținem (cu $K$, compactul mărginit de $\Gamma$):
\[
  \int_\Gamma \vec{V} \cdot d\vec{r} = \iint_K -ydxdy = -\int_{-1}^1 \int_{x^2 - 1}^{\sqrt{1-x^2}} ydxdy.
\]

\vspace{1cm}

10. Să se calculeze fluxul cîmpului $\vec{V} = (x + z^2, y + z^2, -2z)$
prin suprafața
\[
  \Sigma: x^2 + y^2 + z^2 = 1, z \geq 0.
\]

\emph{Indicații:} \emph{Suprafața nu este închisă!} Așadar, nu se poate aplica formula Gauss-Ostrogradski.
Trebuie să o închidem cu cercul $C: x^2 + y^2 = 1, z = 0$, formînd suprafața $S$. Rezultă
$\Sigma \cup C = S$, deci $\kal{F}_\Sigma(\vec{V}) + \kal{F}_C(\vec{V}) = \kal{F}_S(\vec{V})$.

Pe $S$ avem flux nul, deoarece $\nabla \cdot \vec{V} = 0$, iar pe $C$ se poate calcula cu definiția
(vezi și exercițiul 6).

\vspace{1cm}

11. Să se calculeze, folosind formula lui Stokes, integrala curbilinie $\int_\Gamma \alpha$,
pentru:
\[
  \alpha = (y - z)dx + (z - x)dy + (x - y)dz, \quad \Gamma: z = x^2 + y^2, z = 1.
\]

\emph{Indicații:} Definind $\Sigma = \{(x, y, z) \mid z = x^2 + y^2, z \leq 1\}$ avem că
$\Gamma$ este bordul lui $\Sigma$. Putem aplica formula lui Stokes și găsim succesiv:
\begin{align*}
  \int_\Gamma \alpha &= \int_\Sigma -2(dy \wedge dz + dz \wedge dx + dx \wedge dy) \\
                     &= -2 \iint_D (-2x - 2y + 1) dxdy,
\end{align*}
unde $D$ este discul unitate. La primul pas am folosit formula lui Stokes, iar la
al doilea, formula de calcul pentru integrala de suprafață de speța a doua.

Problema se finalizează trecînd la coordonate polare.



\end{document}
